\section{Presentations, interpretations, and classical proof transfer}
\label{sec:logic}

The previous foundation packets already supply the essential numerical
tools: bounded Booleanization, finite witness search, precision-indexed
presentations, and compilation of finite observation contracts. We retain
those tools and add a precise logical interface. The distinction that matters
is between the strength of a mathematical proof and its use of classical
logic. A proof can be classical while using only the nonlogical resources of
primitive recursive arithmetic. For the conclusions specified below, its
classical logical steps can be eliminated without increasing the induction
strength. This is an application of the established negative-translation
method, not a claim of a new conservation theorem; see
\cite{moschovakis} for the broader background. We give the restricted proof
directly, because its preservation of quantifier-free induction is essential.

\subsection{The arithmetic target and its presentations}

Our target \(\ipra\) is the repository's first-order theory
\texttt{PRAAx}. Its variables range over \(\N\). Its function symbols are
finite primitive recursive definitions, represented by \texttt{PR k}; terms
are \texttt{PTm n}; formulas are \texttt{Fm PTm n}; and a completed
formalization must construct a term of type
\(\texttt{Prf PRAAx}\,\Gamma\,\varphi\).
The nonlogical axioms are defining equations, successor separation and
injectivity, and induction for quantifier-free formulas. Intuitionistic
first-order natural deduction supplies the logical rules.

An assertion about a primitive recursive approximation procedure is a
schema: its finite definition is fixed externally and then appears through
the corresponding function symbol. Its numerical parameters, dimensions,
degrees and precisions can all remain quantified. A schema is not a theorem
quantifying over a total universal evaluator for arbitrary primitive recursive
programs. The existing diagonal obstruction to such an evaluator remains in
force.

Finite lists, integers, rationals with positive denominator, polynomial
coefficient arrays and matrix entries have primitive recursive numerical
encodings. A fixed bounded formula about such data has a primitive recursive
Boolean characteristic function. Its equivalence to the characteristic
equation is proved by finite folds and quantifier-free induction, as in the
earlier framework. We use this established bounded reflection principle
throughout. It licenses a bounded invariant; it does not license every true
universal assertion about a primitive recursive function.

\begin{definition}[A numerical presentation of a sort]
\label{def:presented-sort}
A presentation of a mathematical sort consists of an arithmetic domain
predicate \(D(a)\), an arithmetic relation \(E(a,b)\) serving as equality,
and a specified code \(a_0\) with a proof of \(D(a_0)\). On \(D\), the
relation \(E\) must be provably reflexive, symmetric and transitive.
Every operation is represented by a fixed primitive recursive numerical
function, with proofs that it preserves the domains and respects \(E\).
Every additional relation is an arithmetic formula, with a proof that it
respects \(E\) in each argument.
\end{definition}

The domain and equality formulas need not be decidable. For example, for a
fixed family \(a(u,k)\) of rational observations, the assertion that \(u\)
presents a regular name and the equality of two such names are generally
universal arithmetic formulas. This is a specified family of names, not a
code space with a universal evaluator for all primitive recursive functions.
Whenever closure under an operation on name parameters is claimed, the
corresponding numerical constructor and its contracts must be supplied.
Finite-observation compilation is one way of constructing that closure.

\begin{proposition}[Interpretation through numerical presentations]
\label{prop:setoid-interpretation}
Consider a fixed finite many-sorted first-order signature with presentations
as in Definition~\ref{def:presented-sort}. Replace source equality by \(E\),
source operations and relations by their numerical interpretations, and
sorted quantifiers by
\[
 (\forall x{:}D\ A)^I=\forall x\,(D(x)\to A^I),
 \qquad
 (\exists x{:}D\ A)^I=\exists x\,(D(x)\land A^I).
\]
If the interpreted source axioms have \(\ipra\) derivations, every
intuitionistic first-order derivation from those axioms has an interpreted
\(\ipra\) derivation, with domain assumptions for its free variables.
\end{proposition}

\begin{proof}
First induct externally on terms. The closure proofs show that an interpreted
term has the appropriate domain, and the congruence proofs show that replacing
its free parameters by equivalent codes produces an equivalent result.
Then induct externally on formulas to show that the interpreted formula is
unchanged in truth under those replacements. For atomic equality this follows
from symmetry and transitivity of \(E\); for other atoms it is the supplied
relation congruence. The propositional cases follow by the intuitionistic
rules. In a quantified case keep the newly quantified code fixed, use the
induction hypothesis under its domain assumption, and introduce the
restricted quantifier.

Finally transform the given proof by induction on its finite derivation
tree. The propositional rules translate directly. Universal elimination
uses the term-domain lemma to discharge its added domain premise.
Existential introduction supplies both the term-domain proof and the
translated matrix. Existential elimination opens a witness together with
its domain proof; its eigenvariable condition is unchanged. The equality
substitution rule is the formula congruence just established. Source axioms
use their supplied interpretations. The chosen inhabitants ensure the usual
nonempty interpretation of every sort. These are external inductions on
finite syntax and proofs; they add no numerical induction schema to
\(\ipra\).
\end{proof}

This proposition gives an organization of definitions and proofs, rather
than an automatic presentation of every mathematical structure. In
particular, an induction schema in a source theory is an additional axiom
whose translation must be verified. A formula that is quantifier-free in
abstract notation may cease to be quantifier-free after interpretation.
If abstract equality means \(\forall k\,E_k(a,b)\), induction on that
equality is induction on a quantified arithmetic formula. The proposition
does not authorize it. The finite observation method must first replace it
by a numerical invariant at a fixed requested precision.

\subsection{A source theory with classical logic}

Let \(\pra\) denote the theory with exactly the same first-order language
and nonlogical axioms as \(\ipra\), but classical logic. Concretely, add to
the existing intuitionistic deduction system the axiom schema
\[
 A\lor\neg A,
 \qquad \neg A:=(A\to\bot),
\]
for arbitrary first-order formulas \(A\). Thus the source theory is not
Peano arithmetic, a higher-type theory, Lean's ambient mathematics, or set
theory. A proof in a stronger system must first be reduced to the stated
\(\pra\) before the result below can apply.

We use the elementary \(\ipra\) theorem that numerical equality is
decidable, hence stable:
\[
 (s=t)\lor\neg(s=t),
 \qquad \neg\neg(s=t)\longrightarrow s=t.
\]
The repository has an explicit arithmetic route to this fact. In
\texttt{TeqReflect.lean}, the term \texttt{Tm.eqt} reflects scalar equality
by \texttt{Q.eqt\_eq\_zero\_iff}. The cover/selector cases decide whether
the test vanishes. The interpretation \texttt{ToPRA.interp} and its round
trip transport the result to numerical \texttt{PTm} equality. This is an
object-theory proof obligation; a Lean instance of \texttt{DecidableEq}
cannot be substituted for it.

\begin{definition}[Negative translation]
\label{def:negative}
Define \(A\mapsto A^N\) externally by recursion on formulas:
\[
\begin{array}{rcl@{\qquad}rcl}
 (s=t)^N&:=&s=t,& \bot^N&:=&\bot,\\
 (A\land B)^N&:=&A^N\land B^N,
   &(A\to B)^N&:=&A^N\to B^N,\\
 (\forall x A)^N&:=&\forall x A^N,
   &(A\lor B)^N&:=&\neg(\neg A^N\land\neg B^N),\\
 (\exists x A)^N&:=&\neg\forall x\neg A^N.&&&
\end{array}
\]
This atom-preserving form uses the stability of numerical equality.
\end{definition}

Terms and numerical free variables are unchanged. Structural induction on
the formula gives
\[
 (A[\sigma])^N=A^N[\sigma]
\]
for simultaneous capture-avoiding substitution, and the corresponding
weakening identity. At a quantifier both sides use the same lifted
substitution. Moreover, \(A^N\) is quantifier-free whenever \(A\) is
quantifier-free: translating a disjunction introduces only conjunction,
implication and bottom.

\begin{lemma}[Stability of translated formulas]
\label{lem:negative-stable}
For every fixed formula \(A\), \(\ipra\) proves
\(\neg\neg A^N\to A^N\).
\end{lemma}

\begin{proof}
This is external induction on the syntax of \(A\). Equality is stable by
the arithmetic fact above. Bottom is stable by applying a hypothesis
\((\bot\to\bot)\to\bot\) to the identity implication.

Stability is preserved by conjunction: a denial of either conjunct denies
the conjunction, so double negation of the conjunction supplies double
negation of both conjuncts. Apply their stability proofs.
If \(B\) is stable, then \(C\to B\) is stable. Indeed, given
\(\neg\neg(C\to B)\) and \(C\), a denial of \(B\) would deny
\(C\to B\). We therefore obtain \(\neg\neg B\), and then \(B\).
Universal quantification preserves stability by the same argument at an
arbitrary numerical value, followed by universal introduction.

Finally every negation is stable intuitionistically. Given
\(\neg\neg\neg C\), assume \(C\); this implies \(\neg\neg C\),
contradicting the given triple negation. Both the translated disjunction
and translated existential are negations. This proves their stability
without any additional decidability assumption.
\end{proof}

\begin{theorem}[Classical-to-intuitionistic transfer at fixed strength]
\label{thm:negative-transfer}
From a \(\pra\) derivation of \(A\) with finite hypothesis context
\(\Gamma\), one can construct an \(\ipra\) derivation
\[
 \Gamma^N\vdash_{\ipra} A^N.
\]
Here \(\Gamma^N\) translates every context member. The statement requires
a source derivation; truth of the source statement alone does not suffice.
\end{theorem}

\begin{proof}
Transform the source derivation by induction on its finite proof tree.
Conjunction, implication, universal quantification and bottom use the
corresponding intuitionistic deduction rules. Substitution commutation
handles universal instantiation.

For disjunction introduction, a proof of \(A^N\) contradicts the first
conjunct of \(\neg A^N\land\neg B^N\); the other introduction is
symmetric. For disjunction elimination, the translated premises are
\[
 \neg(\neg A^N\land\neg B^N),\qquad
 A^N\to C^N,\qquad B^N\to C^N.
\]
To obtain \(C^N\), use Lemma~\ref{lem:negative-stable}. A denial of
\(C^N\) supplies both \(\neg A^N\) and \(\neg B^N\), contradicting
the first premise.

For existential introduction, an instance \(A^N(t)\) contradicts
\(\forall x\neg A^N(x)\). For existential elimination, the translated
existence premise is \(\neg\forall x\neg A^N(x)\), and the transformed
subderivation obtains \(C^N\) from \(A^N(x)\), with \(x\) fresh for
the conclusion and the other hypotheses. If \(\neg C^N\), this gives
\(\neg A^N(x)\) for arbitrary \(x\); universal introduction contradicts
the existence premise. Stability then gives \(C^N\).

The translation of excluded middle is
\[
 \neg\bigl(\neg A^N\land\neg(A^N\to\bot)\bigr).
\]
Its two prohibited conjuncts would be a formula and its negation, so this
is an intuitionistic theorem. Equality reflexivity is unchanged. For the
equality-substitution rule, all premises asserting equality of numerical
terms are unchanged, and the same rule substitutes equals in the translated
formula.

It remains to check the nonlogical axioms, which is the strength-sensitive
part of the argument. The primitive recursive defining equations,
successor separation and successor injectivity are unchanged. By the
substitution identity, the translation of the induction instance for
\(\varphi\) is precisely
\[
 \varphi^N(0)\to
 \bigl(\forall n(\varphi^N(n)\to\varphi^N(n+1))\bigr)
 \to\forall n\varphi^N(n).
\]
Since \(\varphi\) was quantifier-free, \(\varphi^N\) is also
quantifier-free. This is an allowed \(\ipra\) induction instance.
Every derivation constructor and nonlogical axiom is now covered.
\end{proof}

The transformation is syntactic. With standard finite encodings, it is a
primitive recursive operation on formula and derivation codes: each
constructor inserts a fixed proof pattern, and atomic stability instances
are substitutions into one fixed arithmetic proof. No arbitrary primitive
recursive function symbol is evaluated. A future internal theorem asserting
that this operation sends accepted \(\pra\) proofs to accepted \(\ipra\)
proofs would concern two syntactic proof checkers. It would not assert the
truth of all theorems of either system.

\subsection{Conclusions retained in their original form}

\begin{definition}[A class of retained formulas]
\label{def:retained}
Let \(\mathcal N_b\) be the class generated from fixed bounded formulas
over primitive recursive predicates by conjunction, implication and
unbounded universal quantification. We do not close this class under
unbounded existential quantification or arbitrary disjunction.
\end{definition}

\begin{lemma}[Equivalence to the translation]
\label{lem:negative-retained}
If \(A\in\mathcal N_b\), then
\(\ipra\vdash A^N\leftrightarrow A\).
\end{lemma}

\begin{proof}
First consider a bounded formula. Its bounded subformulas are decidable by
bounded Booleanization. For decidable subformulas, the negative reading of
disjunction is equivalent to ordinary disjunction. For a bounded existential,
scan its finite numerical range. If a candidate succeeds, it supplies the
existential; if all candidates fail, the resulting finite universal failure
contradicts its negative reading. This proves the equivalence for a bounded
existential without unbounded search or Markov's principle. The bounded
universal case uses pointwise equivalence. Equivalently, both readings are
equivalent to the same characteristic Boolean equation.

The remaining closure cases follow by substituting provably equivalent
subformulas in a conjunction, implication or universal quantifier. These
are external inductions on fixed formulas, not object-theory induction on
arbitrary quantified predicates.
\end{proof}

\begin{proposition}[Transfer without changing the asserted conclusion]
\label{prop:negative-conservation}
If every member of \(\Gamma\) and \(A\) belongs to \(\mathcal N_b\), then
\[
 \Gamma\vdash_{\pra} A
 \quad\Longleftrightarrow\quad
 \Gamma\vdash_{\ipra} A.
\]
\end{proposition}

\begin{proof}
Apply Theorem~\ref{thm:negative-transfer}, replace the hypotheses by their
provably equivalent translations, and then use
Lemma~\ref{lem:negative-retained} for the conclusion. The reverse direction
regards an intuitionistic proof as a classical one.
\end{proof}

In particular, a \(\pra\) proof of
\(\forall\mathbf x\,B(\mathbf x)\), with \(B\) bounded, transfers to
\(\ipra\) with the original meaning. Its intermediate proof may use
excluded middle even for unbounded formulas; the target proof contains no
new excluded-middle axiom. This applies to the logical form of exact
polynomial identities, universal finite-list bounds, arithmetic inequalities
and fixed-exponent nonexistence statements. The premise that there is a
\(\pra\) derivation remains essential: the proposition does not turn an
arbitrary published proof or an ambient Lean theorem into such a derivation.

\subsection{Witnesses recovered from bounded classical proofs}

\begin{theorem}[Bounded classical existence yields a verified selector]
\label{thm:bounded-classical-selector}
Fix a primitive recursive relation \(R(\mathbf x,y)\), a primitive
recursive numerical bound \(b(\mathbf x)\), and
\(H(\mathbf x)\in\mathcal N_b\). If \(\pra\) proves
\[
 \forall\mathbf x\left(H(\mathbf x)\to
       \exists y\le b(\mathbf x)\ R(\mathbf x,y)\right),
\]
then an explicitly defined primitive recursive function \(w\) satisfies
\[
 \ipra\vdash\forall\mathbf x\left(H(\mathbf x)\to
       w(\mathbf x)\le b(\mathbf x)\ \land\
       R(\mathbf x,w(\mathbf x))\right).
\]
\end{theorem}

\begin{proof}
The bounded existential belongs to \(\mathcal N_b\); hence the entire
displayed implication does too. Proposition~\ref{prop:negative-conservation}
first gives an \(\ipra\) proof of that implication.

Define \(w\) by scanning \(0,\ldots,b(\mathbf x)\), retaining the first
successful candidate, and returning zero if no candidate succeeds. A
separate found flag distinguishes failure from success at zero. This is
numerical primitive recursion with a given finite bound.

After testing the first \(j\) candidates, use the following bounded
invariant. Either the flag is unset and every tested candidate fails, or
the saved candidate is less than \(j\), satisfies \(R\), and every earlier
candidate fails. Bounded Booleanization makes the invariant
quantifier-free. Its base case is immediate, and its successor step is a
case split on the flag and one primitive recursive test. Quantifier-free
induction therefore proves it. The bounded existential supplied by the
transferred theorem rules out the all-failed terminal state. The saved
candidate is the required witness within its bound.
\end{proof}

This theorem enlarges the acceptable proof interface. One may prove, using
classical \(\pra\) reasoning, that a finite search space contains a desired
certificate, and then obtain the constructive certificate by the fixed
bounded selector. It is unnecessary to redesign every intermediate
classical choice in that existence proof. The substantive requirements are
the primitive recursive search bound and a proof using only the stated
nonlogical strength.

The hypotheses can still describe infinite observations. For example,
regularity of a fixed family of Cauchy names or a uniform local error
contract often has the form \(\forall\mathbf u\,C(\mathbf x,\mathbf u)=1\).
Such a formula belongs to \(\mathcal N_b\), although its validity is not a
decidable numerical test. The theorem's output and search relation remain
finite. The proofs of the existence and accuracy of an operation, and the
data required to execute it, must remain distinct.

\subsection{The boundary and the Lean completion path}

The transfer removes a logical obstacle; it supplies neither missing
quantitative estimates nor missing arithmetic induction reductions. For
example, a finiteness assertion in the form
\[
 \exists H\ \forall z\,
   (\operatorname{Point}(z)\to\operatorname{ht}(z)\le H)
\]
contains an unbounded existential. The theorem translates it to
\[
 \neg\forall H\neg\forall z\,
   (\operatorname{Point}(z)\to\operatorname{ht}(z)\le H),
\]
which does not provide \(H\). An explicit primitive recursive bound is
exactly what makes Theorem~\ref{thm:bounded-classical-selector} stronger
for its restricted conclusion. No general rule
\(\neg\neg\exists y\,R(y)\to\exists y\,R(y)\) is being assumed.
Likewise, stability of equality between represented names does not make
that equality decidable, and does not supply an apartness witness for
division or a singularity classifier for a polynomial branch.

The order of work is consequential. First interpret mathematical objects
and assertions numerically. Then replace every impermissible induction by
a verified finite invariant, fixing observation precision where necessary.
After this arithmetic audit, Theorem~\ref{thm:negative-transfer} can remove
classical logic from a \(\pra\) derivation. Finally a bounded existential
can be realized by Theorem~\ref{thm:bounded-classical-selector}.
Performing negative translation before the arithmetic audit does not cure
induction on an unbounded observation formula.

The companion \texttt{NegativeSyntax.lean} completes the syntax and
axiom-translation layer. It compiles with Lean 4.32.0; its checked declarations
use neither \texttt{sorryAx} nor \texttt{Classical.choice}. The audit records
Lean's standard \texttt{propext} and \texttt{Quot.sound} dependencies in the
metatheory. The completed and remaining steps are concrete:
\begin{enumerate}
\item The file defines \texttt{negative} on the existing \texttt{Fm PTm n},
  and \texttt{ClassicalPRAAx} with one constructor for a \texttt{PRAAx}
  instance and one for excluded middle. It proves the syntax identities
  \texttt{negative\_subst}, \texttt{negative\_shift} and
  \texttt{negative\_qf} by structural induction.
\item Produce the object proof of numerical equality stability through the
  arithmetic route described above. Then construct an actual
  \texttt{Prf PRAAx} proof of stability for every translated formula, using
  the proof patterns in Lemma~\ref{lem:negative-stable}.
\item The file also translates every arithmetic and excluded-middle axiom
  constructor into an actual \texttt{Prf PRAAx} derivation. For induction,
  it invokes \texttt{PRAAx.ind} on \texttt{negative phi};
  \texttt{negative\_qf} certifies the precise side condition. The excluded
  middle translation is proved directly by intuitionistic contradiction.
\item The remaining main task is to prove by induction on the existing
  \texttt{Prf} derivation the target
\[
 \texttt{Prf ClassicalPRAAx}\ \Gamma\ \varphi
 \ \longrightarrow\
 \texttt{Prf PRAAx}\ (\Gamma\texttt{.map negative})\
       (\texttt{negative}\ \varphi).
\]
  The \texttt{eqSubst} case uses the same constructor on the translated
  formula; the quantifier cases use the substitution and weakening
  identities. The critical disjunction and existential elimination cases
  are exactly those written out in Theorem~\ref{thm:negative-transfer}.
\item The file defines retained formulas with quantifier-free leaves.
  Their equivalence and transfer theorems remain named goal propositions,
  as does the full proof translator. The bounded-leaf extension then uses a
  separately proved Booleanization adapter. Derive transfer of retained
  sequents using
  \texttt{Prf.trans} and the equivalence proofs.
\item Implement first-success search using \texttt{PR} symbols and prove its
  bounded found/no-success invariant. The object recurrence rule developed
  in this report supplies the corresponding induction step. This produces
  the selector theorem at the actual \texttt{Prf PRAAx} level.
\end{enumerate}

The human proof of the logical transfer and the selector theorem is complete
above. Their full Lean proof translator and equality-stability adapter are
specified obligations, not completed formalizations. The independently
compiled object recurrence result reported later should not be confused
with completion of this entire translation pipeline.
